
donc @@M0@@~; or @@M1@@
@@M2@@

ii) Par substitution, il suffit de montrer la formule (6.5.3) pour @@M3@@ et pour @@M4@@, @@M5@@. Comme @@M6@@ est sans torsion, on peut plonger cet anneau dans son tensorisé par @@M7@@, et on voit ainsi que @@M8@@ est injectif. Posons @@M9@@, @@M10@@ étant le premier terme non nul et @@M11@@ de @@M12@@. La multiplicativité de @@M13@@ nous donne, d'après (6.5.2) :
@@M14@@
d'où la valeur de @@M15@@ en identifiant les termes de plus bas degré.

6.6. Le @@M16@@-anneau @@M17@@ est augmenté vers @@M18@@~: on définit en effet un @@M19@@-homomorphisme de @@M20@@ dans @@M21@@ par @@M22@@. Il est donc muni d'une filtration canonique de @@M23@@-anneau augmenté. Nous allons la comparer à celle qui vient d'être définie.

\textbf{Proposition 6.6.1.} \uline{Soit} @@M24@@ \uline{un élément de la forme} @@M25@@. \uline{Alors} :
@@M26@@
\uline{les termes non écrits} \uline{étant de degré} @@M27@@.

Par fonctorialité, il suffit de faire la démonstration dans le cas 
